\SourceLeaf{431}%
\[
\begin{aligned}
\frac1r&=A\cos\varphi^{2}+2B\cos\varphi\sin\varphi+C\sin\varphi^{2}\\
&=\frac{A+C}{2}+\frac{A-C}{2}\cos2\varphi+B\sin2\varphi.
\end{aligned}
\]
Setzt man
\[
\begin{aligned}
\frac{A-C}{2}&=E\cos2\theta\\
B&=E\sin2\theta,
\end{aligned}
\]
wobei man annehmen kann, dass $E$ dasselbe Vorzeichen wie $\frac{A-C}{2}$ hat, so ergibt sich
\[
\frac1r=\tfrac12(A+C)+E\cos2(\varphi-\theta).
\]
Offenbar bezeichnet $\varphi$ den Winkel zwischen der Schnittebene und einer anderen Ebene, die durch die Normale und diejenige Tangente verläuft, der die Richtung $M$ entspricht. Also hat $\frac1r$ offenbar seinen größten (absoluten) Wert oder $r$ seinen kleinsten Wert, wenn $\varphi=\theta$ gilt. $\frac1r$ hat seinen kleinsten Wert, wenn $\varphi=\theta+90^{0}$ gilt. Die größte und die kleinste Krümmung treten also in zwei Ebenen auf, die sich im rechten Winkel schneiden. Die Extremwerte für $\frac1r$ sind folglich
\[
\tfrac12(A+C)\pm\BareRoot{\left\{\left(\frac{A-C}{2}\right)^{2}+BB\right\}},
\]
ihre Summe ist $=A+C$ und ihr Produkt $=AC-BB$, oder das Produkt der beiden extremen Krümmungsradien ist
\[
=\frac{1}{AC-BB}.
\]
Dieses Produkt ist von großer Bedeutung und verdient es, genauer untersucht zu werden. Aus den obigen Formeln erhalten wir nämlich
\[
AC-BB=(\alpha\beta'-\beta\alpha')^{2}(TV-UU)ZZ.
\]
Aus der dritten Formel in [Lehrsatz] $6$, Art. $7$, folgt jedoch leicht, dass
\[
\alpha\beta'-\beta\alpha'=\pm Z,
\]
also
\[
AC-BB=Z^{4}(TV-UU).
\]
Ferner gilt nach Art. $8$
